\subsection{Comparison with the Provable Lechowicz Algorithm}
\label{sec:evaluation:lechowicz}

A reviewer concern (and a natural one) is that our headline comparison
is against three heuristics (FIFO, Wait-Awhile, Spatial) rather than
against a provably-optimal online algorithm. We address this by
implementing a FaaS-adapted port of the one-way trading algorithm of
\citet{lechowicz2024online}, which provides an optimal competitive
ratio $\sqrt{U/L}$ for deadline-constrained carbon-aware scheduling
in a simpler model. Our port: for each invocation, examine the carbon
forecast over the SLA-deadline window, compute the threshold $\Phi^*
= \sqrt{U \cdot L}$ (where $U, L$ are the maximum and minimum
forecast intensities), commit at the first slot below $\Phi^*$,
otherwise commit at the deadline.

\begin{table}[h]
\centering\small
\caption{Lechowicz et al.\ \citep{lechowicz2024online} one-way trading
algorithm vs.\ our baselines, 5-seed mean $\pm$ std.}
\label{tab:lechowicz}
\begin{tabular}{llrr}
\toprule
Setup & Scheduler & vs FIFO & Cold-start \\
\midrule
\multirow{4}{*}{Synthetic 5-region}
& Wait-Awhile  & $-306.5\% \pm 4.3\%$ & $2.03\%$ \\
& Lechowicz    & $-238.4\% \pm 3.2\%$ & $2.39\%$ \\
& Spatial      & $+72.9\% \pm 0.3\%$  & $0.03\%$ \\
& \GreenFaaS   & $+71.8\% \pm 0.3\%$  & $0.05\%$ \\
\midrule
\multirow{4}{*}{Real LWA 4-region}
& Wait-Awhile  & $-176.2\% \pm 2.5\%$ & $1.29\%$ \\
& Lechowicz    & $\mathbf{-336.6\% \pm 4.3\%}$ & $2.84\%$ \\
& Spatial      & $+65.1\% \pm 0.2\%$  & $0.03\%$ \\
& \GreenFaaS   & $+64.6\% \pm 0.2\%$  & $0.03\%$ \\
\midrule
\multirow{4}{*}{Real single-region DE}
& Wait-Awhile  & $-176.7\% \pm 6.3\%$ & $1.61\%$ \\
& Lechowicz    & $\mathbf{-320.2\% \pm 2.9\%}$ & $2.44\%$ \\
& Spatial      & $\phantom{+}0.0\% \pm 0.0\%$ & $0.02\%$ \\
& \GreenFaaS   & $\phantom{+}0.0\% \pm 0.0\%$ & $0.02\%$ \\
\bottomrule
\end{tabular}
\end{table}

The finding is striking: \textbf{the provably-optimal-competitive-ratio
algorithm performs worse than the heuristic Wait-Awhile on real
carbon data, and worst of all on the single-region setup where its
theory applies cleanest.} (On \emph{synthetic} carbon the two batch
schedulers rank the other way --- Lechowicz $-238.4\%$ vs
Wait-Awhile $-306.5\%$, so Lechowicz is marginally the less
catastrophic --- but on real carbon the ordering reverses, Lechowicz
$-336.6\%$ vs Wait-Awhile $-176.2\%$; either way both fail badly.)
In the most-favourable-to-Lechowicz setup
(real single-region DE, no spatial axis), \GreenFaaS{} and Spatial
correctly identify zero opportunity and match FIFO byte-for-byte,
while Lechowicz inflates carbon by 320\% above FIFO.

The mechanism is exactly the
\S\ref{sec:tradeoff:idle} contribution. Lechowicz's competitive
analysis assumes deferral is free of side effects: the cost of holding
a job pending until commit time is treated as zero. For FaaS, this
assumption is violated by warm-pool idle energy: deferring an
invocation extends the warm container's idle period, charging
$P_i \Delta t \bar{C}_r$ in carbon that the competitive analysis does
not see. The threshold $\Phi^* = \sqrt{U \cdot L}$ aggressively
defers any invocation arriving above $\Phi^*$, even when the warm
container's idle-energy cost will exceed the execution savings ---
which is, given typical FaaS values of $\beta = P_i T_w / (P_a
\tau_e) \in [50, 200]$, almost always. The cold-start rate of
2.4--2.8\% (vs.\ \GreenFaaS's 0.02--0.05\%) is the direct evidence: the
algorithm aggressively defers past warm-pool TTLs and pays cold-start
penalties on most invocations.

This is, in our reading, the strongest empirical evidence for the
paper's central technical claim. A provably-optimal algorithm for
deadline-constrained one-way trading, when ported to FaaS scheduling
without accounting for warm-pool idle-energy coupling, becomes the
worst-performing scheduler in our study. The closed-form idle-energy
correction of \S\ref{sec:tradeoff:idle} is precisely what makes the
difference; without it, provability does not rescue the algorithm.
\S\ref{sec:theory} formalizes the warm-pool-coupled cost model
and states Conjecture~\ref{conj:lech-lb}: that
Lechowicz-style algorithms have $\Omega(\beta)$ competitive ratio
in a multi-container regime that captures FaaS production behavior.
We have not proven this conjecture; the single-container model does
not admit the natural N-invocation construction (the idle energy of
one shared container cannot be charged per pending invocation without
double-counting, see \S\ref{sec:theory:why_hard}), and the empirical
evidence here is consistent with the conjecture but does not
establish it.

\paragraph{Honest caveats.} (i) Our port follows Lechowicz et al.'s
algorithmic prescription faithfully but does not extend their
competitive-ratio proof to the warm-pool-coupled cost. The non-separable
cost structure of \GreenFaaS{} (\S\ref{sec:problem:online}) is
outside the standard one-way trading framework. (ii) The Lechowicz
framework was designed for batch jobs with hours-long deadlines,
similar to Wait-Awhile's intended regime; we are not surprised that a
FaaS-adapted port fails, and the paper's finding is that
\emph{neither} batch-oriented framework (heuristic or provably-
optimal) transfers to FaaS without addressing the idle-energy
coupling. A FaaS-native online algorithm with competitive-ratio
guarantees over the warm-pool-coupled problem remains the most
interesting theoretical extension this work suggests.
(iii) Our port computes the one-way trading bounds $U$ and $L$ from
the carbon-intensity forecast over each invocation's deadline window,
whereas the original algorithm assumes globally-known a-priori
bounds. Using deadline-local bounds adapts the threshold
$\Phi^* = \sqrt{UL}$ to each invocation's forecast window; this is a
necessary adaptation to the non-stationary carbon setting, and it may
make the threshold more or less aggressive than the global-bound
version depending on the local carbon range.
